% Generated draft: requires mathematical and source review.
\documentclass[11pt]{article}
\usepackage[margin=25mm]{geometry}
\usepackage{fontspec}
\setmainfont{Latin Modern Roman}
\usepackage{amsmath,amssymb,mathtools,unicode-math}
\setmathfont{Latin Modern Math}
\usepackage{enumitem,needspace,xurl,hyperref}
\setlength{\parindent}{0pt}
\setlength{\parskip}{5pt}
\setlength{\emergencystretch}{3em}
\newcommand{\N}{\mathbb{N}}
\newcommand{\Z}{\mathbb{Z}}
\newcommand{\Q}{\mathbb{Q}}
\newcommand{\R}{\mathbb{R}}
\newcommand{\C}{\mathbb{C}}
\newcommand{\F}{\mathbb{F}}
\newcommand{\A}{\mathbb{A}}
\newcommand{\PP}{\mathbb{P}}
\newcommand{\E}{\mathbb{E}}
\newcommand{\eps}{\varepsilon}
\newcommand{\dd}{\,\mathrm{d}}
\begin{document}
\section*{UnsolvedMath ID: 30005848}
\textbf{Source problem code:} catalogue entry
\subsection{Definitions and mathematical statement}
For an oriented compact Riemannian manifold $(M,g)$ of dimension $n$, write $d_g$ for its Riemannian distance, $T_M$ for its current of integration, and $\operatorname{Vol}_g(M)$ for its volume. The Sormani--Wenger intrinsic-flat distance between two such integral current spaces is
\[
 d_{\mathcal F}(M_1,M_2)
 =\inf\bigl\{\mathbf M(A)+\mathbf M(B):
 \varphi_{1\#}T_{M_1}-\varphi_{2\#}T_{M_2}
 =A+\partial B\bigr\},
\]
where the infimum ranges over complete metric spaces $Z$, distance-preserving embeddings $\varphi_i:(M_i,d_{g_i})\to Z$, and integral currents $A$ and $B$ in $Z$ of dimensions $n$ and $n+1$, respectively; $\mathbf M$ denotes current mass. A flat $n$-torus is $\R^n/\Lambda$ with its induced Euclidean metric, where $\Lambda\subset\R^n$ is a full-rank lattice.

\textbf{Requested formulation (incomplete in the supplied source).} Formulate and prove a Geroch stability theorem with hypotheses that exclude bubbling and drawstring degeneration, and with a volume-preserving intrinsic-flat convergence conclusion. In notation, the desired conclusion for the sequence covered by such a theorem is convergence to some flat torus $(T^n,g_\infty)$ with
\[
 d_{\mathcal F}\bigl((M_j,d_{g_j},T_{M_j}),
 (T^n,d_{g_\infty},T_{T^n})\bigr)\longrightarrow0,
 \qquad
 \operatorname{Vol}_{g_j}(M_j)
 \longrightarrow\operatorname{Vol}_{g_\infty}(T^n).
\]
These equations specify the requested type of conclusion, not a theorem asserted for an identified class of sequences. The dimension $n$, topology and boundary conditions of the approximating manifolds, orientations, scalar-curvature assumptions, normalization and other geometric bounds, precise exclusion conditions, and whether passage to a subsequence is permitted are not specified in the supplied excerpt.

The original Question 1 says that an entropy bound in its Theorem 2 excludes splines and bubbles along a sequence, then asks for a Geroch stability formulation excluding bubbles and drawstrings. Theorem 2 and its entropy bound are not supplied. The accompanying figure description characterizes splines as becoming arbitrarily thin, a bubble neck as pinching to a point, and a drawstring as pulling a circle almost to a point, with that circle collapsing to a point along the sequence. These descriptions do not supply quantitative definitions of the excluded degenerations. Consequently, a fully quantified stability conjecture cannot be recovered without additional source information or an explicitly proposed formulation.
\subsection{Short English statement}
The problem asks for hypotheses preventing bubbles and drawstrings under which a Geroch stability theorem yields convergence to a flat torus in Sormani--Wenger intrinsic-flat distance, together with convergence of total volumes. The supplied excerpt leaves the admissible sequences and exclusion hypotheses unspecified, so this is a request to formulate a theorem rather than a complete conjecture. The supplied literature status is \emph{partially solved}; it is dated curation requiring independent review.
\subsection{Sources}
\begin{itemize}
\item Local UnsolvedMath registry, permanent ID \texttt{30005848}, title \emph{Geroch Stability Without Bubbles and Drawstrings}, source problem code \texttt{catalogue entry}, in the Oberwolfach Reports Open Problems collection. No individual proposer is identified in the supplied record; the proposed year is 2024.
\item \emph{Analysis, Geometry and Topology of Positive Scalar Curvature Metrics} (2024), Oberwolfach Reports, pp.\ 483--566, DOI \texttt{10.4171/owr/2024/9}. The supplied original excerpt identifies Question 1 and an accompanying Figure 1 on p.\ 491. \url{https://doi.org/10.4171/owr/2024/9}.
\item \emph{Stability of the positive mass theorem and torus rigidity}. Supplied literature reference: \url{https://arxiv.org/abs/2210.04340}.
\item \emph{Conjectures on convergence and scalar curvature}. Supplied literature reference: \url{https://arxiv.org/abs/2103.10093}.
\item Literature assessment checked at \texttt{2026-08-22}: this is supplied dated curation, not independently verified here. It labels the problem partially solved and reports that intrinsic-flat stability results cover important torus and positive-mass regimes under curvature, diameter, or integral bounds, while the proposed bubble-free Geroch stability statement in its full generality remains unproved. The supplied background expressly requests independent verification of time-sensitive status and cited claims.
\item The registry labels the statement \texttt{corrected\_verified} and says that it was reconstructed and checked against the official Oberwolfach report. These are supplied provenance claims, not independent source verification performed for this draft.
\end{itemize}
\end{document}
